% !PRINTED 257
\weberpage{257}

If now $n$ is greater than 1, then one can take $\varrho$ and $1+\varrho$ to be different from zero, and then set $\sigma=-1:(1+\varrho)$, whereby the substitution (9) goes over into $A_0$.

But if $n=1$ and $p=2$, then the \gesperrt{first exceptional case} occurs; for then one of the two numbers $\varrho$ and $1+\varrho$ is always equal to zero. This case leads to a group of degree $6$
\[
E = \begin{pmatrix}1,&0\\0,&1\end{pmatrix}, \begin{pmatrix}1,&0\\1,&1\end{pmatrix}, \begin{pmatrix}0,&1\\1,&1\end{pmatrix}, \begin{pmatrix}1,&1\\1,&0\end{pmatrix}, \begin{pmatrix}1,&1\\0,&1\end{pmatrix}, \begin{pmatrix}0,&1\\1,&0\end{pmatrix},
\]
in which a normal divisor of degree $3$ is contained:
\[
G = \begin{pmatrix}1,&0\\0,&1\end{pmatrix}, \begin{pmatrix}0,&1\\1,&1\end{pmatrix}, \begin{pmatrix}1,&1\\1,&0\end{pmatrix}.
\]

We now set this exceptional case aside. We then have further:

4) The group $G$ contains every substitution of the form
\weq{10}{\begin{pmatrix} 1, & 0 \\ \xi, & 1 \end{pmatrix},}
where $\xi$ is an arbitrary number in $\mathfrak{E}$, except in the case $n=1$, $p=3$.

For by 3), $G$ contains the substitution
\[
\begin{pmatrix} 0, & \varrho \\ -\varrho^{-1}, & 0 \end{pmatrix}\begin{pmatrix} 0, & 1 \\ -1, & 0 \end{pmatrix}\begin{pmatrix} 0, & -\varrho \\ \varrho^{-1}, & 0 \end{pmatrix}\begin{pmatrix} 0, & 1 \\ -1, & 0 \end{pmatrix} = \begin{pmatrix} \varrho^2, & 0 \\ 0, & \varrho^{-2} \end{pmatrix},
\]
if $\varrho$ is an arbitrary number in $\mathfrak{E}$ different from zero. From this it follows that
\weq{11}{\begin{pmatrix} 1, & 0 \\ -\lambda, & 1 \end{pmatrix}\begin{pmatrix} \varrho^{-2}, & 0 \\ 0, & \varrho^2 \end{pmatrix}\begin{pmatrix} 1, & 0 \\ \lambda, & 1 \end{pmatrix}\begin{pmatrix} \varrho^2, & 0 \\ 0, & \varrho^{-2} \end{pmatrix} = \begin{pmatrix} 1, & 0 \\ \lambda(\varrho^4-1), & 1 \end{pmatrix}}
also occurs in $G$ for every arbitrary $\lambda$. If now one can take $\varrho$ different from zero in such a way that $\varrho^4-1$ does not vanish, then for every $\xi$ one can determine the number $\lambda$ from $\lambda(\varrho^4-1)=\xi$, and thus obtains from (11) every substitution of the form (10). In order to judge the possibility of this, we take a primitive root $\gamma$ of $\mathfrak{E}$ and set $\varrho=\gamma^h$.

If one can take the exponent $h$ so that $4h$ is not divisible by $p^n-1$, then $\varrho$ satisfies our requirement. This is always possible when $p=2$, $n>1$, and when $p^n-1>4$, and so only the two cases $n=1$, $p=3$ and $n=1$, $p=5$ still remain doubtful.
% [signature line "Weber, Algebra. II.  17" omitted -- gathering signature, not body text]
